% =============================================================================
%  VLM4RWD @ NeurIPS 2026 -- workshop submission (8 pages excl. refs/appendix)
%  Double-blind: keep anonymous.
%
%  NOTE: swap the preamble for the official neurips_2026.sty when available:
%      \usepackage[final]{neurips_2026}   % camera-ready
%      \usepackage{neurips_2026}          % anonymous submission
%  The article-based preamble below compiles standalone in the meantime.
% =============================================================================
\documentclass[10pt]{article}
\usepackage[margin=1in]{geometry}
\usepackage{amsmath,amssymb}
\usepackage{booktabs}
\usepackage{graphicx}
\usepackage{xcolor}
\usepackage{multirow}
\usepackage[hidelinks]{hyperref}
\usepackage{natbib}

% --- visible TODO markers; set to \relax before submitting ---
\newcommand{\todo}[1]{\textcolor{red}{\textbf{[TODO: #1]}}}
\newcommand{\pending}{\textcolor{red}{--}}

\title{Forgetting Without Hallucinating:\\
Calibrated Concept Unlearning in Vision--Language Models}
\author{Anonymous Submission}
\date{}

\begin{document}
\maketitle

% =============================================================================
\begin{abstract}
Concept unlearning methods for vision--language models are evaluated almost
entirely by how far accuracy on the erased concept falls. We show that this
number hides \emph{how} the model fails. A standard gradient-ascent unlearner
drives accuracy on the target concept to zero across all domains, but does so by
\emph{confidently} assigning those inputs to a single semantic neighbour: its
predictions on forgotten inputs have an effective support of roughly
$4.5$ classes out of $126$. The erasure succeeds and the model hallucinates, and
the standard metric cannot distinguish the two.
We attribute this to \emph{unbounded} suppression: gradient ascent on the target
class drives the image feature into a degenerate region of the representation,
after which the freed probability mass collapses onto whichever class is
nearest. We propose a minimal correction --- suppression bounded by a cap,
combined with a redistribution term --- together with a retain-set construction
that allows forgetting supervised in a single domain to propagate to domains the
forgetting loss never observes. On Mini-DomainNet with CLIP ViT-B/16, our method
matches gradient ascent on erasure ($0\%$ target accuracy in all four domains)
and on retention ($85.2\%$), while the model's predictions on forgotten inputs
remain calibrated: mean confidence $22.1\%$ and an effective support of $\sim
37$ classes.
\todo{add detectability AUROC sentence once measured}
\end{abstract}

% =============================================================================
\section{Introduction}
\label{sec:intro}

Vision--language models such as CLIP are increasingly deployed as general-purpose
recognisers in settings where their behaviour must be revisable after training.
Regulatory requirements, licensing disputes, and safety constraints can all
demand that a specific concept be removed from a deployed model, and retraining
from scratch is rarely practical. This has motivated a body of work on
\emph{concept unlearning} for vision--language models, including the
domain-agnostic setting in which a concept must be erased not only where it was
observed but across every visual domain the model encounters. We adopt that
setting rather than propose it; our concern is not with how erasure is achieved
but with what the model does \emph{afterwards}.

Almost universally, unlearning is judged by a single quantity: accuracy on the
erased concept, reported as it falls toward zero, alongside accuracy on the
retained classes. This pair of numbers answers whether the concept is still
recognised and whether the rest of the model survived. It is silent on a third
question that matters at deployment time: when the unlearned model is shown an
instance of the erased concept, \emph{what does it predict instead, and with what
confidence?} A model that has genuinely lost a concept and a model that has
merely been redirected to a confident substitute produce identical forget
accuracy.

We find that this distinction is not hypothetical. Applying a standard
gradient-ascent unlearning objective to CLIP ViT-B/16 on Mini-DomainNet erases
the target concept completely --- accuracy falls to $0\%$ in all four domains ---
while the model's predictions on those same inputs remain sharply peaked, with
an effective support of roughly $4.5$ classes out of $126$ and the majority of
the freed probability mass landing on a single semantic neighbour. The unlearned
model does not abstain; it confidently asserts a specific wrong label. By the
reported metric the procedure is a complete success, and by the model's own
output it is a fluent, high-confidence error --- precisely the kind of
unsupported assertion that undermines reliability in deployment.

We trace this behaviour to the unbounded nature of the suppression signal.
Gradient ascent on the target class drives the image feature to be increasingly
anti-aligned with the target text embedding, and CLIP's learned temperature
amplifies the effect; in our runs the objective diverges and the representation
degrades. Once the representation is degraded, the probability freed by
suppression has nowhere to go except onto whichever class remains nearest, which
is what produces the confident substitute. This suggests that the choice is not
between forgetting and calibration, but that calibration is lost as a
\emph{side effect} of how aggressively forgetting is applied.

Acting on this, we study a minimal correction: the same cross-entropy ascent,
but \emph{bounded} by a cap that halts the push once the target's probability
falls below chance, combined with terms that encourage the freed mass to spread
both within and across inputs. We pair this with a retain-set construction in
which the target concept is withheld from every domain, which is what allows
forgetting supervised in one domain to reach domains the forgetting loss never
observes. On Mini-DomainNet the resulting model matches gradient ascent on both
quantities the literature reports --- $0\%$ target accuracy in all four domains,
$85.2\%$ retain accuracy --- while its predictions on forgotten inputs remain
uncertain: mean confidence $22.1\%$ and an effective support of $\sim\!37$
classes. The two methods are indistinguishable under the standard evaluation and
behave very differently under inspection.

We view the practical consequence as the more important one. A model that is
uncertain on erased concepts exposes those inputs to ordinary confidence-based
monitoring; a model that is confidently wrong routes them silently into
downstream decisions. Reporting only forget accuracy cannot distinguish these
two deployments.

\paragraph{Contributions.}
\begin{itemize}
  \item We show that a standard unlearning objective achieves complete
        cross-domain erasure while producing \emph{confidently incorrect}
        predictions on the erased concept --- a failure mode that forget
        accuracy, the field's primary metric, cannot detect.
  \item We attribute this to unbounded suppression degrading the representation,
        and show that bounding the suppression preserves the model's capacity to
        express uncertainty at equal erasure and retention.
        \todo{requires the cap sweep, Table~\ref{tab:ablation}(b)}
  \item We report, alongside forget and retain accuracy, the confidence and
        effective support of predictions on erased inputs, and argue these should
        accompany any unlearning result.
  \item We evaluate over \todo{$N$} concepts and \todo{$S$} seeds, stratified by
        semantic neighbourhood structure, and include a retain-only control
        isolating the contribution of the suppression term.
\end{itemize}

% =============================================================================
\section{Related Work}
\label{sec:related}

\paragraph{Concept and domain unlearning in CLIP.}
\todo{Open with prior work honestly and explicitly. Cross-domain / domain-agnostic
concept removal has been posed before; we do not claim the problem setting. Cite
and describe: null-space projection for selective and domain-agnostic unlearning;
CLIPErase; zero-shot CLIP class forgetting via text--image space adaptation
(TMLR 2025) and via synthetic samples (WACV 2025); null-space calibration (UNSC).
State plainly what is different here: we study the \emph{output behaviour} of
unlearned models rather than proposing a new erasure mechanism.}

\paragraph{Approximate unlearning.}
\todo{gradient ascent / NegGrad lineage; SISA and exact retraining; note that
evaluation in this literature is dominated by forget/retain accuracy and MIA.}

\paragraph{Calibration and failure-mode analysis.}
\todo{output entropy and confidence as calibration measures; effective support
size (perplexity). Be explicit that these are standard quantities, not ours.}

% =============================================================================
\section{Setup}
\label{sec:setup}

\paragraph{Model.} CLIP ViT-B/16 with the pretrained backbone frozen. We train
only the prompt-tuning components: $8$ learnable context tokens on each of the
text and vision branches inserted into the first $9$ layers, bottleneck
adapters, and an instance-wise prompt generator (cross-attention at layer $9$).
No pretrained CLIP weight is modified.

\paragraph{Data.} Mini-DomainNet: $126$ classes over $4$ domains (clipart,
painting, real, sketch), $8$-shot per (class, domain) for the retain set.

\paragraph{Task.} Given a target concept observed in a \emph{single} domain
(e.g.\ \texttt{tiger} in \emph{sketch}), the model must fail to recognise that
concept in \emph{all four} domains while the remaining $125$ classes are
preserved in every domain.

\paragraph{Evaluation set.} All metrics are computed on held-out images: every
train-split image not used for training, with the few-shot set and the forget
pool excluded by path. \todo{state final counts}

% =============================================================================
\section{Method}
\label{sec:method}

Each step draws from two streams: the $8$-shot retain loader, and a
\emph{forget pool} of up to $100$ images of the target concept from the forget
domain (a random subset of $24$ is forwarded per step, which preserves pool
diversity while bounding memory and giving an unbiased per-step estimate of the
marginal term).

The objective is $\mathcal{L} = \mathcal{L}_{\text{ret}} + \lambda_f
\mathcal{L}_{\text{forget}}$.

\paragraph{Retain term.} Cross-entropy on the retain batch, with the target
class \emph{excluded from every domain}. The model is trained on the other $125$
classes in all four domains and receives no positive signal for the target
anywhere.

\paragraph{Forget term.} On forget images, with $\ell$ the scaled logits and $t$
the target index:
\begin{align}
\mathcal{L}_{\text{sup}} &= \min\!\big(\mathrm{CE}(\ell, t),\; \tau\big)
  && \text{(maximised; \emph{bounded} suppression)} \label{eq:sup}\\
\mathcal{H}_{\text{cond}} &= \tfrac{1}{m}\textstyle\sum_i H\!\big(p_i^{\neg t}\big)
  && \text{(maximised; per-image spread)} \label{eq:hcond}\\
\mathcal{H}_{\text{marg}} &= H\!\big(\tfrac{1}{m}\textstyle\sum_i p_i^{\neg t}\big)
  && \text{(maximised; across-image spread)} \label{eq:hmarg}
\end{align}
so that $\mathcal{L}_{\text{forget}} = -\mathcal{L}_{\text{sup}}
- \lambda_c \mathcal{H}_{\text{cond}} - \lambda_m \mathcal{H}_{\text{marg}}$,
where $p_i^{\neg t}$ is the softmax over non-target classes.

\paragraph{Why the cap.} Equation~\eqref{eq:sup} is gradient ascent on the target
cross-entropy --- the standard forgetting signal --- with one modification. Left
unbounded, ascent drives the image feature to be maximally anti-aligned with the
target text embedding; CLIP's learned temperature ($\approx\!100$) amplifies
this, and in our runs the loss diverged to $-195$ and degraded the model. The cap
$\tau=6$ halts the push once the target's probability is $\approx 0.25\%$, below
chance ($1/126 \approx 0.79\%$), leaving the representation intact. Our claim is
that this is precisely what preserves the model's capacity to spread the freed
probability rather than collapsing it.

\paragraph{Why the target is excluded from retain everywhere.} This is what lets
forgetting cross domains. Retaining the target in the untouched domains supplies
a restoring gradient and confines forgetting to the supervised domain; removing
it does not. All domains are scored against the same text embeddings and share
the same prompt parameters, so suppression driven by one domain's images
modifies parameters every domain depends on. Table~\ref{tab:ablation} isolates
this.

% =============================================================================
\section{Experiments}
\label{sec:exp}

\paragraph{Protocol.} \todo{Final protocol: $N$ concepts $\times$ $S$ seeds,
identical grid for every method. Concepts stratified by neighbour structure
(close neighbour / mid / isolated) to test whether the redistribution behaviour
depends on semantic proximity. Report mean $\pm$ std.}

% -----------------------------------------------------------------------------
\begin{table}[t]
\centering
\caption{Main comparison. Target concept \texttt{tiger}, forget domain
\emph{sketch}, seed 1, single run. \textbf{All numbers below are from a single
concept and a single seed and are placeholders for the full grid.}
``Eff.\ support'' is $\exp(H)$ of the prediction on forgotten inputs (perplexity;
$126$ = uniform). Higher is more calibrated.}
\label{tab:main}
\small
\begin{tabular}{lcccccc}
\toprule
& \multicolumn{2}{c}{Target acc.\ (\%) $\downarrow$}
& \multicolumn{2}{c}{Preserved (\%) $\uparrow$}
& \multicolumn{2}{c}{Behaviour on forgotten inputs} \\
\cmidrule(lr){2-3}\cmidrule(lr){4-5}\cmidrule(lr){6-7}
Method & forget dom. & other doms. & retain & neighbour
& conf.\ (\%) $\downarrow$ & eff.\ support $\uparrow$ \\
\midrule
Original (no unlearning)      & 94.9 & \pending & \pending & \pending & \pending & \pending \\
Retain-only (control)         & \pending & \pending & \pending & \pending & \pending & \pending \\
Entropy max.\ (ADU-style)     & 53.1 & 91.7 & 85.7 & 91.7 & \pending & \pending \\
NegGrad + retain              & \textbf{0.0} & 0.3 & $\sim$85 & 93.3 & high & 4.5 \\
\midrule
Ours                          & \textbf{0.0} & \textbf{0.0} & 85.2 & 92.0
                              & \textbf{22.1} & \textbf{37.0} \\
\bottomrule
\end{tabular}
\end{table}
% -----------------------------------------------------------------------------

\paragraph{Reading Table~\ref{tab:main}.} NegGrad and our method are
indistinguishable on the metric the literature reports: both erase the concept in
every domain, both preserve the retain set. They differ entirely in
\emph{how} the model behaves on the erased concept --- an effective support of
$4.5$ versus $37$ classes.

\todo{Retain-only row is the control that shows suppression, not exclusion,
causes the $0\%$. Must be filled before any claim in \S\ref{sec:method} stands.}

% -----------------------------------------------------------------------------
\begin{table}[t]
\centering
\caption{Ablations. \todo{single-seed placeholders; rerun on the full grid.}}
\label{tab:ablation}
\small
\begin{tabular}{lccccc}
\toprule
Variant & target (forget dom.) & target (other doms.) & retain
& conf.\ (\%) & eff.\ support \\
\midrule
\multicolumn{6}{l}{\emph{(a) Retain-set construction}}\\
\quad target kept in retain (anchor on) & 29.5 & 88.0 & 86.0 & 79.5 & 2.5 \\
\quad target excluded (ours)            & 0.0  & 0.0  & 85.2 & 22.1 & 37.0 \\
\midrule
\multicolumn{6}{l}{\emph{(b) Suppression cap $\tau$}}\\
\quad $\tau = \infty$ (unbounded)       & \pending & \pending & \pending & \pending & \pending \\
\quad $\tau = 8$                        & \pending & \pending & \pending & \pending & \pending \\
\quad $\tau = 6$ (ours)                 & 0.0 & 0.0 & 85.2 & 22.1 & 37.0 \\
\midrule
\multicolumn{6}{l}{\emph{(c) Forget-pool size}}\\
\quad $8$ (few-shot only)               & 0.0 & 0.0 & 85.5 & 57.2 & 6.4 \\
\quad $100$ (ours)                      & 0.0 & 0.0 & 85.2 & 22.1 & 37.0 \\
\midrule
\multicolumn{6}{l}{\emph{(d) Redistribution term}}\\
\quad variance-min instead of entropy   & 0.0 & 0.0 & 85.7 & 30.2 & 22.1 \\
\quad per-image entropy only            & 0.0 & 0.0 & 85.5 & 29.1 & 23.2 \\
\quad + batch-marginal (ours)           & 0.0 & 0.0 & 85.2 & 22.1 & 37.0 \\
\bottomrule
\end{tabular}
\end{table}
% -----------------------------------------------------------------------------

\paragraph{The cap is the load-bearing component.} \todo{This paragraph is the
paper's central mechanistic claim and is currently \emph{unsupported}. Block (b)
of Table~\ref{tab:ablation} must be run. Expected: $\tau=\infty$ diverges and
collapses effective support; bounded $\tau$ preserves it at equal erasure.}

\paragraph{Detectability.} \todo{Because the model is uncertain on forgotten
inputs, a confidence threshold should separate them from retained inputs. Report
AUROC for ours vs.\ NegGrad. This is the deployment consequence and the strongest
answer to ``why does calibration matter here''. Not yet run.}

% =============================================================================
\section{Analysis}
\label{sec:analysis}

\paragraph{Where the freed probability goes.}
\todo{Report the redistribution honestly, against the runner-up reference rather
than against uniform. Zero-shot CLIP with the target logit masked --- what a
model without the concept predicts --- concentrates $35.1\%$ on a single class
(effective support $\sim 6$); ours concentrates $30.2\%$ (effective
support $\sim 21$). Conclusion to state: residual concentration reflects genuine
semantic similarity in the data, not a pathology introduced by the method.
Note also that isotropic features would give $3.4\%$, so this is \emph{not} an
architectural limit of CLIP --- do not claim that.}

\paragraph{Erasure is shallow.}
\todo{Subtracting a single fixed bias vector from the logits restores the target
class as the top-1 prediction for $46.5\%$ of forgotten inputs. Currently
measured only on our own model with an oracle-estimated offset. Before this
appears as a claim it needs (i) held-out estimation of the offset and (ii) the
same probe applied to NegGrad. If it holds on a method we did not design, it is a
statement about unlearning; otherwise it is a limitation of ours and belongs in
\S\ref{sec:limits}.}

\paragraph{Negative results.}
\todo{Appendix material. (i) Feature effective dimensionality does not predict
redistribution concentration ($3.6$ dims yields a lower concentration than
$12.3$). (ii) The first two moments of the forget-logit distribution do not
determine it (TV $=0.218$ against a matched Gaussian). (iii) Flattening the mean
logit vector makes concentration worse, not better. These rule out the obvious
mechanistic explanations.}

% =============================================================================
\section{Limitations}
\label{sec:limits}
\begin{itemize}
  \item Single dataset (Mini-DomainNet) and a single backbone (ViT-B/16).
  \item We do not compare against null-space projection; no code is available and
        we judged an unfaithful reimplementation worse than none.
  \item The suppression term is gradient ascent on cross-entropy, the same family
        as NegGrad; the bound, not the ascent, is what differs.
  \item Driving accuracy \emph{below} chance is itself a detectable signature of
        deliberate erasure.
  \item \todo{if the recovery probe holds: our erasure is shallow and recoverable.}
\end{itemize}

% =============================================================================
\section{Conclusion}
\todo{Short. Forget accuracy alone is not evidence of good unlearning; the
behaviour of the model on forgotten inputs is a separate and reportable axis.}

% =============================================================================
\bibliographystyle{plainnat}
\bibliography{refs}

\appendix
\section{Additional results}
\todo{negative results; sensitivity to $\lambda_c$ (non-monotone --- larger
weights increase concentration); per-mode cluster analysis; runner-up reference
table; compute cost (ours requires training and forget data, unlike training-free
projection baselines --- state this plainly).}

\end{document}
